\documentclass{ximera}
\title{Differentials}
\begin{document}
\begin{abstract}
Treating $\frac{dy}{dx}$ as a fraction to define the differentials $dx$ and $dy$, and using $dy$ as an approximation of $\Delta y$. Textbook §3.10.
\end{abstract}
\maketitle

Remember how we had said that $\frac{d y}{d x}$ is not a fraction and we shouldn't think of it as one? We're about to change all of that!

% CHECKED-FIX: was "We call $d x$ and $y$" (prose typo, not an \answer)
We can think of $d y$ as a very small change, an approximation of a larger $\Delta y$. Let's just pretend that $f^{\prime}(x)=\frac{d y}{d x}$ is a fraction and put like terms on the same side. So we have $dy=\answer{f^{\prime}(x)\,dx}$. We call $d x$ and $d y$ \wordChoice{\choice[correct]{differentials}\choice{derivatives}\choice{increments}\choice{limits}}.

It turns out we can use differentials as a good approximation of $\Delta y$. Let's see an example to see what in the world I'm talking about.

\begin{example}\label{example:9-1}
Let $f(x)=x^{2}-2 x+4$. Let us look at the change from 2 to 2.01.
% work: f(2)=4-4+4=4, f(2.01)=4.0401-4.02+4=4.0201, so Delta y = 0.0201
\[
\Delta y=f(2.01)-f(2)=\answer{0.0201}
\]
% work: f'(x)=2x-2, dx = 2.01-2 = 0.01, dy = f'(2) dx = 2(0.01) = 0.02
\[
f^{\prime}(x)=\answer{2x-2}, \qquad dy=f^{\prime}(2)\,dx=\answer{0.02}
\]
\end{example}

So this ``division'' turns out to be a good approximation! Also, finding $d y$ was a lot easier to compute than $\Delta y$.

\end{document}
